
\documentclass{article}
\usepackage{colortbl}
\usepackage{makecell}
\usepackage{multirow}
\usepackage{supertabular}

\begin{document}

\newcounter{utterance}

\centering \large Interaction Transcript for game `mmlu\_pro', experiment `mmlu\_pro\_biology', episode 374 with Qwen3.6{-}27B.
\vspace{24pt}

{ \footnotesize  \setcounter{utterance}{1}
\setlength{\tabcolsep}{0pt}
\begin{supertabular}{c@{$\;$}|p{.15\linewidth}@{}p{.15\linewidth}p{.15\linewidth}p{.15\linewidth}p{.15\linewidth}p{.15\linewidth}}
    \# & \multicolumn{2}{c}{Player} && \multicolumn{2}{c}{Game Master} \\
    \hline

    \theutterance \stepcounter{utterance}  
    & & & \multicolumn{4}{p{0.6\linewidth}}{
        \cellcolor[rgb]{0.9,0.9,0.9}{
            \makecell[{{p{\linewidth}}}]{
                \texttt{\tiny{[P1$\langle$GM]}}
                \texttt{The following are multiple choice questions (with answers) about biology. Answer only with "The answer is (X)" where X is the correct letter choice.} \\
\\ 
\texttt{Question:} \\
\texttt{Which of the following represents an accurate statement concerning arthropods?} \\
\texttt{Options:} \\
\texttt{A. They possess an exoskeleton composed primarily of peptidoglycan.} \\
\texttt{B. They possess an open circulatory system with a dorsal heart.} \\
\texttt{C. They are members of a biologically unsuccessful phylum incapable of exploiting diverse habitats and nutrition sources.} \\
\texttt{D. They lack paired, jointed appendages.} \\
\texttt{E. N/A} \\
\texttt{F. N/A} \\
\texttt{G. N/A} \\
\texttt{H. N/A} \\
\texttt{I. N/A} \\
\texttt{J. N/A} \\
\texttt{The answer is (B)} \\
\\ 
\texttt{Question:} \\
\texttt{In a given population, 1 out of every 400 people has a cancer caused by a completely recessive allele, b. Assuming the population is in Hardy{-}Weinberg equilibrium, which of the following is the expected proportion of individuals who carry the b allele but are not expected to develop the cancer?} \\
\texttt{Options:} \\
\texttt{A. 19/400} \\
\texttt{B. 1/400} \\
\texttt{C. 40/400} \\
\texttt{D. 38/400} \\
\texttt{E. 2/400} \\
\texttt{F. 1/200} \\
\texttt{G. 20/400} \\
\texttt{H. 50/400} \\
\texttt{I. N/A} \\
\texttt{J. N/A} \\
\texttt{The answer is (D)} \\
\\ 
\texttt{Question:} \\
\texttt{A mutation in a bacterial enzyme changed a previously polar amino acid into a nonpolar amino acid. This amino acid was located at a site distant from the enzyme’s active site. How might this mutation alter the enzyme’s substrate specificity?} \\
\texttt{Options:} \\
\texttt{A. By changing the enzyme’s pH optimum} \\
\texttt{B. By changing the enzyme's molecular weight} \\
\texttt{C. An amino acid change away from the active site increases the enzyme's substrate specificity.} \\
\texttt{D. By changing the shape of the protein} \\
\texttt{E. By changing the enzyme's temperature optimum} \\
\texttt{F. By altering the enzyme's ability to be denatured} \\
\texttt{G. By changing the enzyme’s location in the cell} \\
\texttt{H. By changing the enzyme's color} \\
\texttt{I. An amino acid change away from the active site cannot alter the enzyme’s substrate specificity.} \\
\texttt{J. By altering the enzyme's rate of reaction} \\
\texttt{The answer is (D)} \\
\\ 
\texttt{Question:} \\
\texttt{Which of the following is not a way to form recombinant DNA?} \\
\texttt{Options:} \\
\texttt{A. Translation} \\
\texttt{B. Conjugation} \\
\texttt{C. Specialized transduction} \\
\texttt{D. Transformation} \\
\texttt{E. N/A} \\
\texttt{F. N/A} \\
\texttt{G. N/A} \\
\texttt{H. N/A} \\
\texttt{I. N/A} \\
\texttt{J. N/A} \\
\texttt{The answer is (A)} \\
\\ 
\texttt{Question:} \\
\texttt{Which of the following is not known to be involved in the control of cell division?} \\
\texttt{Options:} \\
\texttt{A. Microtubules} \\
\texttt{B. Checkpoints} \\
\texttt{C. DNA polymerase} \\
\texttt{D. Centrosomes} \\
\texttt{E. Cyclins} \\
\texttt{F. Mitochondria} \\
\texttt{G. Protein kinases} \\
\texttt{H. Fibroblast cells} \\
\texttt{I. N/A} \\
\texttt{J. N/A} \\
\texttt{The answer is (H)} \\
\\ 
\texttt{Question:} \\
\texttt{In human, the M, MN, and N blood groups are determined by twocodominantalleles: L\^{}m andL\^{}n. In a population of 800 college students the following phenotypes were identified: Blood Group Number M 392 MN 336 N 72 800 Determine the frequency of the Lm and I" alleles.} \\
\texttt{Options:} \\
\texttt{A. 0.65 for L\^{}m and 0.35 for L\^{}n} \\
\texttt{B. 0.60 for L\^{}m and 0.40 for L\^{}n} \\
\texttt{C. 0.5 for both L\^{}m and L\^{}n} \\
\texttt{D. 0.8 for L\^{}m and 0.2 for L\^{}n} \\
\texttt{E. 0.75 for L\^{}m and 0.25 for L\^{}n} \\
\texttt{F. 0.9 for L\^{}m and 0.1 for L\^{}n} \\
\texttt{G. 0.7 for L\^{}m and 0.3 for L\^{}n} \\
\texttt{H. 0.85 for L\^{}m and 0.15 for L\^{}n} \\
\texttt{I. 0.55 for L\^{}m and 0.45 for L\^{}n} \\
\texttt{J. 0.6 for L\^{}m and 0.4 for L\^{}n} \\
            }
        }
    }
     \\ \\

    \theutterance \stepcounter{utterance}  
    & \multicolumn{4}{p{0.6\linewidth}}{
        \cellcolor[rgb]{0.561,0.737,0.561}{
            \makecell[{{p{\linewidth}}}]{
                \texttt{\tiny{[P1$\rangle$GM]}}
                \texttt{The answer is (A)} \\
            }
        }
    }
    & & \\ \\

    \theutterance \stepcounter{utterance}  
    & & & \multicolumn{2}{p{0.3\linewidth}}{
        \cellcolor[rgb]{0.9,0.9,0.9}{
            \makecell[{{p{\linewidth}}}]{
                \texttt{\tiny{[GM$|$GM]}}
                \texttt{A} \\
            }
        }
    }
    & & \\ \\

    \theutterance \stepcounter{utterance}  
    & & & \multicolumn{2}{p{0.3\linewidth}}{
        \cellcolor[rgb]{0.9,0.9,0.9}{
            \makecell[{{p{\linewidth}}}]{
                \texttt{\tiny{[GM$|$GM]}}
                \texttt{G} \\
            }
        }
    }
    & & \\ \\

    \theutterance \stepcounter{utterance}  
    & & & \multicolumn{2}{p{0.3\linewidth}}{
        \cellcolor[rgb]{0.9,0.9,0.9}{
            \makecell[{{p{\linewidth}}}]{
                \texttt{\tiny{[GM$|$GM]}}
                \texttt{game\_result = LOSE} \\
            }
        }
    }
    & & \\ \\

\end{supertabular}
}

\end{document}
